% ============================================================
%  Chapitre 1 : Les ensembles
% ============================================================
\chapter{Les ensembles}
\chaptermark{Les ensembles}

\begin{quelqueslignes}
L'ensemble est la brique élémentaire de toutes les mathématiques appliquées à l'informatique : les tables d'une base de données, les fichiers d'un répertoire, les étudiants d'une promotion ou les adresses d'un réseau local se décrivent tous comme des collections d'objets bien définis. Ce chapitre fixe le vocabulaire (appartenance, représentation, cardinal), puis construit les trois opérations fondamentales (union, intersection, complémentaire) et leurs propriétés, avec le diagramme de Venn comme support visuel. Il prépare directement le chapitre 2, où les mêmes idées deviennent des propositions logiques, et le chapitre 3, où compter les codages possibles revient à compter les sous-ensembles d'un ensemble fini.
\end{quelqueslignes}

\section{Notions d'ensembles}

\begin{definition}[ensemble, appartenance]
Un \emph{ensemble} est une collection d'objets, appelés \emph{éléments}, bien définie : pour tout objet $x$, on peut dire sans ambiguïté si $x$ appartient à l'ensemble (on écrit $x \in A$) ou n'y appartient pas (on écrit $x \notin A$). Un ensemble peut se noter en extension, en écrivant tous ses éléments entre accolades, ou en compréhension, en donnant la propriété que ses éléments vérifient.
\end{definition}

Les deux écritures suivantes désignent le même ensemble des entiers de 3 à 9 :
\[
A = \{3, 4, 5, 6, 7, 8, 9\} \qquad\text{et}\qquad A = \{x \in \N \mid 3 \leqslant x \leqslant 9\}.
\]
L'écriture en extension est directe mais ne convient qu'aux ensembles finis et assez courts ; l'écriture en compréhension est la seule possible pour un ensemble infini comme $\N$ lui-même. L'ordre des éléments dans les accolades n'a aucune importance, et écrire deux fois le même élément ne change rien : $\{1, 2, 2, 3\}$ est exactement l'ensemble $\{1, 2, 3\}$.

\begin{definition}[cardinal, ensemble fini]
Un ensemble est \emph{fini} lorsque ses éléments peuvent être comptés. Le nombre de ses éléments s'appelle son \emph{cardinal} et se note $|A|$. Le \emph{vide}, noté $\varnothing$, est l'ensemble qui ne contient aucun élément : $|\varnothing| = 0$. Un ensemble à un seul élément s'appelle un \emph{singleton}.
\end{definition}

\begin{exemple}[représenter et compter]
Soit $A = \{x \in \N \mid 3 \leqslant x \leqslant 9\}$ et $B = \{c, t, r, p\}$, l'ensemble des lettres du mot « cptr ». L'ensemble $A$ s'écrit en extension $\{3, 4, 5, 6, 7, 8, 9\}$ et son cardinal vaut $|A| = 7$ ; l'ensemble $B$ contient quatre lettres distinctes, donc $|B| = 4$. Enfin $5 \in A$, $10 \notin A$ et la lettre $t \in B$, alors que $z \notin B$.
\end{exemple}

Les ensembles de nombres que nous utiliserons pendant tout le semestre forment une chaîne d'emboîtement : chaque ligne du tableau~\ref{ch1:tab:nombres} contient la précédente.

\begin{center}
\small
\begin{tabular}{L{1.6cm}L{7.4cm}L{6.6cm}}
\toprule
\rowcolor{robleuf}\textbf{Notation} & \textbf{Ensemble} & \textbf{Exemples d'éléments}\\
\midrule
$\N$ & entiers naturels : $0, 1, 2, 3, \dots$ & $0$ ; $17$ ; $256$\\
\rowcolor{rogrisf!50}$\Z$ & entiers relatifs : $\dots, -2, -1, 0, 1, 2, \dots$ & $-7$ ; $0$ ; $1024$\\
$\D$ & nombres décimaux, écriture finie après la virgule & $-7$ ; $0{,}75$ ; $2{,}5$\\
\rowcolor{rogrisf!50}$\Q$ & nombres rationnels, quotients d'entiers & $-7$ ; $\dfrac{2}{3}$ ; $0{,}75$\\
$\R$ & nombres réels, rationnels et irrationnels & $\sqrt{2}$ ; $\pi$ ; $0{,}75$\\
\bottomrule
\end{tabular}
\captionof{table}{Les ensembles de nombres usuels, emboîtés de $\N$ vers $\R$}
\label{ch1:tab:nombres}
\end{center}

Le tableau~\ref{ch1:tab:nombres} se lit de gauche à droite comme des poupées gigognes : tout entier naturel est un entier relatif, tout entier relatif est un décimal, tout décimal est un rationnel, tout rationnel est un réel. Les décimaux $\D$ sont les nombres dont l'écriture s'arrête après la virgule, comme $0{,}75$ ou $2{,}5$ : chacun s'écrit comme un quotient d'entiers par une puissance de 10, et la fraction $\tfrac{1}{3}$, dont les décimales se répètent sans fin, n'en fait pas partie. La figure~\ref{ch1:fig:nombres} illustre cet emboîtement. Une petite difficulté de vocabulaire mérite d'être signalée dès maintenant : le nombre $\dfrac{4}{2}$, bien qu'écrit comme un quotient, vaut exactement 2 et appartient donc à $\N$ ; c'est la valeur qui décide de l'appartenance, non l'écriture. Les nombres $\sqrt{2}$, $\pi$ ou $0{,}123456789101112\dots$ sont, eux, irrationnels : leurs décimales ne se répètent jamais de façon périodique, et aucune paire d'entiers ne les représente exactement.

\begin{center}
\begin{tikzpicture}[scale=0.92]
  % --- R : reels (bande du haut : etiquettes et exemples) ---
  \draw[rogris!70, line width=0.8pt, fill=rogrisf!30, rounded corners=2mm]
    (0,0) rectangle (14.2,6.0);
  \node[anchor=west, text=rogris, font=\footnotesize\bfseries] at (0.35,5.5) {$\R$ réels};
  \node[anchor=east, text=rogris, font=\footnotesize] at (13.85,5.5) {$\sqrt{2}$ \quad $\pi$};
  % --- Q : rationnels ---
  \draw[robleu, line width=0.8pt, fill=robleuf, rounded corners=2mm]
    (0.55,0.35) rectangle (12.2,5.0);
  \node[anchor=west, text=robleu, font=\footnotesize\bfseries] at (0.9,4.55) {$\Q$ rationnels};
  \node[anchor=east, text=rogris, font=\footnotesize] at (11.85,4.55) {$\dfrac{2}{3}$};
  % --- D : decimaux ---
  \draw[roor, line width=0.8pt, fill=roorf, rounded corners=2mm]
    (1.1,0.7) rectangle (10.2,4.05);
  \node[anchor=west, text=roor, font=\footnotesize\bfseries] at (1.45,3.6) {$\D$ décimaux};
  \node[anchor=east, text=rogris, font=\footnotesize] at (9.85,3.6) {$0{,}75$ \quad $2{,}5$};
  % --- Z : entiers relatifs ---
  \draw[roteal, line width=0.8pt, fill=rotealf, rounded corners=2mm]
    (1.65,1.05) rectangle (8.2,3.15);
  \node[anchor=west, text=roteal, font=\footnotesize\bfseries] at (2.0,2.7) {$\Z$ entiers relatifs};
  \node[anchor=east, text=rogris, font=\footnotesize] at (7.85,2.7) {$-7$};
  % --- N : entiers naturels ---
  \draw[rogold, line width=0.8pt, fill=rogoldf, rounded corners=2mm]
    (2.2,1.4) rectangle (7.0,2.25);
  \node[anchor=west, text=rogold, font=\footnotesize\bfseries] at (2.5,1.82) {$\N$ naturels};
  \node[anchor=east, text=rogris, font=\footnotesize] at (6.7,1.82) {$0$ ; $17$ ; $256$};
\end{tikzpicture}
\captionof{figure}{Les ensembles de nombres s'emboîtent : $\N \subset \Z \subset \D \subset \Q \subset \R$}
\label{ch1:fig:nombres}
\end{center}

\begin{exemple}[classer des nombres]
Rangeons les nombres $-7$, $\dfrac{2}{3}$, $0{,}75$, $0$, $\sqrt{2}$, $15$ dans les bons ensembles. Les nombres $-7$, $0$ et $15$ sont des entiers relatifs, donc ils appartiennent à $\Z$ et aussi à $\N$ pour $0$ et $15$ ; le nombre $0{,}75$ est décimal, car $0{,}75 = \tfrac{75}{100}$ ; le nombre $\dfrac{2}{3}$ est rationnel sans être décimal, ses décimales se répétant sans fin ; le nombre $\sqrt{2}$ est réel mais irrationnel, car aucun quotient d'entiers ne vaut exactement $\sqrt{2}$. On a donc $0, 15 \in \N$ ; $-7 \in \Z$ ; $0{,}75 \in \D$ ; $\dfrac{2}{3} \in \Q$ ; $\sqrt{2} \in \R$ mais $\sqrt{2} \notin \Q$.
\end{exemple}

\section{Diagramme de Venn}

Le diagramme de Venn représente un ensemble par une courbe fermée dessinée à l'intérieur d'un rectangle : le rectangle figure l'\emph{univers} $U$, c'est-à-dire l'ensemble de tous les objets considérés pendant l'exercice, et chaque courbe regroupe les éléments d'un ensemble. Cette représentation visuelle permet de lire les trois opérations fondamentales avant même de les calculer.

\begin{definition}[union, intersection, complémentaire]
Soient $A$ et $B$ deux ensembles inclus dans un univers $U$.
\begin{itemize}
  \item l'\emph{union} $A \cup B$ est l'ensemble des éléments qui appartiennent à $A$ \emph{ou} à $B$ (au moins l'un des deux) ;
  \item l'\emph{intersection} $A \cap B$ est l'ensemble des éléments qui appartiennent à $A$ \emph{et} à $B$ (aux deux à la fois) ;
  \item le \emph{complémentaire} de $A$, noté $\overline{A}$, est l'ensemble des éléments de $U$ qui n'appartiennent pas à $A$.
\end{itemize}
\end{definition}

\begin{center}
\begin{tikzpicture}[scale=0.8]
  \draw[rogris!60, line width=0.6pt, fill=rogrisf!40] (0,0) rectangle (11.2,5.4);
  \node[text=rogris, font=\footnotesize] at (0.55,4.95) {$U$};
  \draw[robleu, line width=1pt, fill=robleuf, fill opacity=0.55] (4.0,2.7) ellipse (2.6 and 1.9);
  \draw[roteal, line width=1pt, fill=rotealf, fill opacity=0.55] (7.2,2.7) ellipse (2.6 and 1.9);
  \node[text=robleu, font=\bfseries] at (2.6,4.1) {$A$};
  \node[text=roteal, font=\bfseries] at (8.6,4.1) {$B$};
  \node[font=\footnotesize, text=rogris] at (3.1,2.7) {$A \cap \overline{B}$};
  \node[font=\footnotesize, text=rogris] at (5.6,2.7) {$A \cap B$};
  \node[font=\footnotesize, text=rogris] at (8.1,2.7) {$\overline{A} \cap B$};
  \node[font=\footnotesize, text=rogris] at (1.75,0.55) {$\overline{A \cup B}$};
  \node[font=\footnotesize, text=rogris] at (9.45,0.55) {$\overline{A \cup B}$};
\end{tikzpicture}
\captionof{figure}{Diagramme de Venn de deux ensembles sécants : chaque région porte un nom, et $\overline{A} \cap \overline{B} = \overline{A \cup B}$}
\label{ch1:fig:venn}
\end{center}

La figure~\ref{ch1:fig:venn} se lit ainsi : l'union $A \cup B$ couvre les deux courbes réunies, l'intersection $A \cap B$ est le lentille central où elles se chevauchent, et le complémentaire de l'union est la zone du rectangle extérieure aux deux courbes. Pour compter les éléments d'un diagramme, on commence toujours par remplir la zone centrale $A \cap B$, puis on en déduit les régions $A$ seul et $B$ seul par soustraction.

\begin{exemple}[répartir 120 étudiants de CPTEC]
Dans une promotion de $120$ étudiants de l'Institut Supérieur CPTEC, $65$ suivent le module de programmation, $50$ suivent le module de réseaux et $25$ suivent les deux. Remplissons le diagramme de Venn, l'univers étant la promotion entière. La zone centrale vaut $|P \cap R| = 25$. La zone « programmation seulement » vaut $65 - 25 = 40$, la zone « réseaux seulement » vaut $50 - 25 = 25$, et l'union couvre donc $40 + 25 + 25 = 90$ étudiants. Il reste $120 - 90 = 30$ étudiants qui ne suivent aucun des deux modules : ils occupent la zone $\overline{P \cup R}$. On vérifie au passage la formule $|P \cup R| = |P| + |R| - |P \cap R| = 65 + 50 - 25 = 90$.
\end{exemple}

\begin{center}
\begin{tikzpicture}[scale=0.8]
  \draw[rogris!60, line width=0.6pt, fill=rogrisf!40] (0,0) rectangle (11.2,5.4);
  \node[text=rogris, font=\footnotesize] at (0.55,4.95) {$U = 120$};
  \draw[robleu, line width=1pt, fill=robleuf, fill opacity=0.55] (4.0,2.7) ellipse (2.6 and 1.9);
  \draw[roteal, line width=1pt, fill=rotealf, fill opacity=0.55] (7.2,2.7) ellipse (2.6 and 1.9);
  \node[text=robleu, font=\bfseries] at (2.6,4.1) {$P$};
  \node[text=roteal, font=\bfseries] at (8.6,4.1) {$R$};
  \node[font=\footnotesize, text=rogris] at (3.1,2.7) {$40$};
  \node[font=\bfseries, text=roor] at (5.6,2.7) {$25$};
  \node[font=\footnotesize, text=rogris] at (8.1,2.7) {$25$};
  \node[font=\footnotesize, text=rogris] at (5.6,0.5) {$30$};
\end{tikzpicture}
\captionof{figure}{Répartition des 120 étudiants : 90 suivent au moins un des deux modules}
\label{ch1:fig:cptec}
\end{center}

\begin{exemple}[compter avec le complémentaire]
Dans une classe de $40$ étudiants, $18$ participent au club de mathématiques ($M$) et $23$ au club d'informatique ($I$), $12$ étudiants participant aux deux clubs. L'union couvre $|M \cup I| = 18 + 23 - 12 = 29$ étudiants ; le complémentaire $\overline{M \cup I}$ compte donc $40 - 29 = 11$ étudiants qui ne participent à aucun club. Les lois de De Morgan le confirment : ne participer à aucun club, c'est ne pas être au club de mathématiques \emph{et} ne pas être au club d'informatique, soit $\overline{M} \cap \overline{I} = \overline{M \cup I}$, et les deux lectures donnent bien les mêmes $11$ étudiants.
\end{exemple}

Les trois opérations se composent : on peut former $\overline{A \cap B}$, $(A \cup B) \cap C$ et bien d'autres. Le diagramme de Venn suit sans effort : avec trois ensembles, le rectangle se partage en huit régions, comme le montre la figure~\ref{ch1:fig:venn3}. Chaque région correspond à une combinaison d'appartenances (dans $A$ mais ni dans $B$ ni dans $C$, dans $A$ et $B$ mais pas dans $C$, et ainsi de suite), et le comptage se fait toujours du centre vers l'extérieur : d'abord le noyau $A \cap B \cap C$, puis les intersections doubles, puis les zones simples.

\begin{center}
\begin{tikzpicture}[scale=0.8]
  \draw[rogris!60, line width=0.6pt, fill=rogrisf!40] (0,0) rectangle (10.4,5.8);
  \node[text=rogris, font=\footnotesize] at (0.5,5.35) {$U$};
  \draw[robleu, line width=1pt, fill=robleuf, fill opacity=0.5] (3.9,3.3) ellipse (2.3 and 1.7);
  \draw[roteal, line width=1pt, fill=rotealf, fill opacity=0.5] (6.5,3.3) ellipse (2.3 and 1.7);
  \draw[rogold, line width=1pt, fill=rogoldf, fill opacity=0.5] (5.2,2.1) ellipse (2.3 and 1.7);
  \node[text=robleu, font=\bfseries] at (2.5,4.55) {$A$};
  \node[text=roteal, font=\bfseries] at (7.9,4.55) {$B$};
  \node[text=rogold, font=\bfseries] at (5.2,0.68) {$C$};
  \node[font=\scriptsize, text=rogris] at (3.3,4.1) {1};
  \node[font=\scriptsize, text=rogris] at (5.2,4.35) {2};
  \node[font=\scriptsize, text=rogris] at (7.1,4.1) {3};
  \node[font=\scriptsize, text=rogris] at (3.6,2.55) {4};
  \node[font=\scriptsize, text=rogris] at (6.8,2.55) {5};
  \node[font=\scriptsize, text=rogris] at (4.0,1.0) {6};
  \node[font=\bfseries, text=rogris] at (5.2,2.9) {7};
  \node[font=\scriptsize, text=rogris] at (1.2,0.9) {8};
\end{tikzpicture}
\captionof{figure}{Trois ensembles sécants : huit régions numérotées, dont le noyau $A \cap B \cap C$ (région 7)}
\label{ch1:fig:venn3}
\end{center}

\section{Lien entre deux ensembles}

Deux ensembles peuvent être comparés de trois manières : l'un peut être contenu dans l'autre, ils peuvent être égaux, ou ils peuvent ne rien partager.

\begin{definition}[inclusion, égalité, ensembles disjoints]
Soient $A$ et $B$ deux ensembles.
\begin{itemize}
  \item $A$ est \emph{inclus} dans $B$, on écrit $A \subset B$, lorsque tout élément de $A$ appartient aussi à $B$ ; on dit que $A$ est un \emph{sous-ensemble} de $B$.
  \item $A$ et $B$ sont \emph{égaux}, on écrit $A = B$, lorsqu'ils ont exactement les mêmes éléments, ce qui revient à avoir à la fois $A \subset B$ et $B \subset A$.
  \item $A$ et $B$ sont \emph{disjoints} lorsque $A \cap B = \varnothing$ : aucun élément commun.
\end{itemize}
\end{definition}

\begin{propriete}[caractérisation de l'inclusion]
Pour tous ensembles $A$ et $B$ inclus dans un univers $U$ :
\[
A \subset B \quad\Longleftrightarrow\quad A \cap B = A \quad\Longleftrightarrow\quad A \cup B = B .
\]
Ces deux égalités sont très utilisées : elles traduisent l'inclusion en une relation entre les opérations des sections précédentes.
\end{propriete}

Par exemple, avec $A = \{1, 2, 3\}$, $B = \{1, 2, 3, 4, 5\}$ et $C = \{6, 7\}$, on a $A \subset B$ car les trois éléments de $A$ figurent dans $B$, et $A \cap C = \varnothing$ : les ensembles $A$ et $C$ sont disjoints. En revanche $B \not\subset A$ car $4 \in B$ alors que $4 \notin A$.

Une confusion classique sépare l'\emph{appartenance} de l'\emph{inclusion} : l'élément 2 appartient à $\{1, 2, 3\}$, on écrit $2 \in \{1, 2, 3\}$, tandis que l'ensemble $\{2\}$ est inclus dans $\{1, 2, 3\}$, on écrit $\{2\} \subset \{1, 2, 3\}$. L'appartenance relie un objet à un ensemble, l'inclusion relie un ensemble à un ensemble ; écrire $2 \subset \{1, 2, 3\}$ ou $\{2\} \in \{1, 2, 3\}$ n'a pas de sens ici, et cette distinction de niveaux évite bien des erreurs de rédaction.

\begin{definition}[produit cartésien]
Le \emph{produit cartésien} de deux ensembles $A$ et $B$, noté $A \times B$, est l'ensemble des couples $(x, y)$ avec $x \in A$ et $y \in B$. Lorsque $A$ et $B$ sont finis, son cardinal vaut $|A \times B| = |A| \times |B|$.
\end{definition}

\begin{exemple}[tester l'inclusion et la disjonction]
Avec $A = \{1, 2, 3\}$, $B = \{1, 2, 3, 4, 5\}$ et $C = \{6, 7\}$, on a $A \subset B$ car les trois éléments de $A$ figurent dans $B$, et l'on vérifie la caractérisation : $A \cap B = \{1, 2, 3\} = A$ et $A \cup B = \{1, 2, 3, 4, 5\} = B$. Les ensembles $A$ et $C$ sont disjoints, puisque $A \cap C = \varnothing$. En revanche $B \not\subset A$ car $4 \in B$ alors que $4 \notin A$ : un seul contre-exemple suffit à ruiner une inclusion.
\end{exemple}

Un couple $(x, y)$ est une paire ordonnée : $(1, 2)$ et $(2, 1)$ sont deux couples distincts. Cette notion décrit toutes les combinaisons possibles de deux choix successifs : un identifiant et un mot de passe, une ligne et une colonne d'une feuille de calcul, un utilisateur et une machine d'un réseau. Ainsi, un service qui propose 3 identifiants possibles et 2 mots de passe autorisés offre $3 \times 2 = 6$ couples identifiant-mot de passe : $|A \times B| = 3 \times 2 = 6$.

\begin{exemple}[compter les sous-ensembles]
Pour construire un sous-ensemble, chaque élément de $A$ offre exactement deux choix : être pris ou être laissé. Un sous-ensemble de $A$ se décrit donc entièrement par un mot de trois bits, qui marque par 1 les éléments pris et par 0 les éléments laissés :

\begin{center}
\small
\begin{tabular}{C{1.9cm}C{1.9cm}C{1.9cm}L{3.2cm}L{2.6cm}}
\toprule
\rowcolor{robleuf}\textbf{$a$ pris ?} & \textbf{$b$ pris ?} & \textbf{$c$ pris ?} & \textbf{Sous-ensemble} & \textbf{Mot binaire}\\
\midrule
0 & 0 & 0 & $\varnothing$ & 000\\
\rowcolor{rogrisf!50}1 & 0 & 0 & $\{a\}$ & 100\\
1 & 1 & 0 & $\{a, b\}$ & 110\\
\rowcolor{rogrisf!50}1 & 1 & 1 & $\{a, b, c\}$ & 111\\
\bottomrule
\end{tabular}
\captionof{table}{Chaque sous-ensemble correspond à un mot binaire}
\label{ch1:tab:bits}
\end{center}

Il y a donc $2 \times 2 \times 2 = 2^{3} = 8$ sous-ensembles, dont quatre sont listés dans le tableau~\ref{ch1:tab:bits} et les quatre autres sont $\{b\}$, $\{c\}$, $\{a,c\}$ et $\{b,c\}$. Le résultat est général : un ensemble à $n$ éléments possède $2^{n}$ sous-ensembles. Ce nombre est exactement celui des codages binaires sur $n$ bits, étudié au chapitre 3 : compter les codages, c'est compter les sous-ensembles.
\end{exemple}

\section{Propriétés des opérations sur les ensembles}

Les opérations $\cup$, $\cap$ et le complémentaire obéissent à des lois fixes, rassemblées dans le tableau~\ref{ch1:tab:lois}. Ces lois se lisent en remplaçant mentalement $A$ et $B$ par des zones du diagramme de Venn ; chacune se traduit par un dessin où les deux membres de l'égalité recouvrent exactement la même région.

\begin{center}
\small
\begin{tabular}{L{4.2cm}L{5.7cm}L{5.7cm}}
\toprule
\rowcolor{robleuf}\textbf{Loi} & \textbf{Première forme} & \textbf{Seconde forme}\\
\midrule
Commutativité & $A \cup B = B \cup A$ & $A \cap B = B \cap A$\\
\rowcolor{rogrisf!50}Associativité & $(A \cup B) \cup C = A \cup (B \cup C)$ & $(A \cap B) \cap C = A \cap (B \cap C)$\\
Distributivité & $A \cap (B \cup C) = (A \cap B) \cup (A \cap C)$ & $A \cup (B \cap C) = (A \cup B) \cap (A \cup C)$\\
\rowcolor{rogrisf!50}Idempotence & $A \cup A = A$ & $A \cap A = A$\\
Absorption & $A \cup (A \cap B) = A$ & $A \cap (A \cup B) = A$\\
\rowcolor{rogrisf!50}Involutivité & \multicolumn{2}{l}{$\overline{\overline{A}} = A$}\\
Lois de De Morgan & $\overline{A \cup B} = \overline{A} \cap \overline{B}$ & $\overline{A \cap B} = \overline{A} \cup \overline{B}$\\
\rowcolor{rogrisf!50}Neutre et absorbant & $A \cup \varnothing = A$, \ $A \cup U = U$ & $A \cap U = A$, \ $A \cap \varnothing = \varnothing$\\
\bottomrule
\end{tabular}
\captionof{table}{Lois algébriques des opérations sur les ensembles}
\label{ch1:tab:lois}
\end{center}

Chaque loi se lit sur le diagramme de Venn en coloriant les deux membres : ils recouvrent exactement la même région. L'absorption $A \cup (A \cap B) = A$ se lit en un coup d'œil : la zone de $A \cap B$ est déjà à l'intérieur de $A$, donc l'union ne l'agrandit pas. L'idempotence $A \cup A = A$ dit simplement qu'unionner un ensemble avec lui-même ne change rien, de même que la distributivité gauche ou droite se vérifie en superposant deux coloriages. Cette lecture graphique ne remplace pas le calcul : elle l'oriente, en suggérant quelle loi appliquer avant de le mener.

Les lois de De Morgan méritent une lecture attentive : le complémentaire d'une union est l'intersection des complémentaires, et le complémentaire d'une intersection est l'union des complémentaires. En français courant : « ne suivre ni programmation ni réseaux » revient à « ne pas suivre programmation \emph{et} ne pas suivre réseaux ». On les retrouve mot pour mot au chapitre 2 sous la forme des lois de De Morgan de la logique.

\begin{exemple}[vérifier une distributivité sur des ensembles concrets]
Prenons $A = \{1, 2, 3, 4\}$, $B = \{3, 4, 5\}$ et $C = \{4, 5, 6\}$ et vérifions $A \cap (B \cup C) = (A \cap B) \cup (A \cap C)$. D'un côté, $B \cup C = \{3, 4, 5, 6\}$, donc $A \cap (B \cup C) = \{3, 4\}$. De l'autre, $A \cap B = \{3, 4\}$ et $A \cap C = \{4\}$, donc $(A \cap B) \cup (A \cap C) = \{3, 4\} \cup \{4\} = \{3, 4\}$. Les deux membres coïncident : la loi est vérifiée sur cet exemple, comme elle l'est toujours.
\end{exemple}

\begin{remarque}[la même algèbre pour les ensembles et pour la logique]
Le tableau~\ref{ch1:tab:lois} et le tableau des connecteurs du chapitre 2 se ressemblent fort : union et intersection jouent les rôles du « ou » et du « et », le complémentaire joue le rôle de la négation, et les lois portent les mêmes noms. Ce n'est pas un hasard : « $x$ appartient à $A \cap B$ » est une phrase qui dit « $x \in A$ \emph{et} $x \in B$ ». Le chapitre 2 exploitera ce parallèle jusqu'au bout, et l'algèbre de Boole le formalisera pour les circuits.
\end{remarque}

\begin{exemple}[vérifier une loi de De Morgan sur des ensembles concrets]
Dans l'univers $U = \{1, 2, 3, 4, 5, 6, 7, 8\}$, prenons $A = \{2, 3, 5\}$ et $B = \{1, 2, 4, 5, 6\}$ et vérifions $\overline{A \cup B} = \overline{A} \cap \overline{B}$. D'un côté, $A \cup B = \{1, 2, 3, 4, 5, 6\}$, donc $\overline{A \cup B} = \{7, 8\}$. De l'autre, $\overline{A} = \{1, 4, 6, 7, 8\}$ et $\overline{B} = \{3, 7, 8\}$, donc $\overline{A} \cap \overline{B} = \{7, 8\}$. Les deux membres coïncident : la loi est vérifiée sur cet exemple, et elle l'est pour tous les autres.
\end{exemple}

\begin{remarque}
Une loi du tableau~\ref{ch1:tab:lois} se retient d'autant mieux qu'on l'a vue échouer une fois dans son usage. Par exemple, la distributivité de $\cup$ sur $\cap$ est souvent oubliée par les étudiants : $A \cup (B \cap C)$ n'est \emph{pas} égal en général à $(A \cup B) \cap C$, mais à $(A \cup B) \cap (A \cup C)$. En cas de doute, refaire un petit test numérique comme celui de l'exemple précédent : trois minutes de calcul évitent une erreur d'orientation.
\end{remarque}

\section{Simplification à l'aide des propriétés}

Beaucoup d'expressions d'ensembles s'écrivent sous une forme bien plus courte. La démarche est algébrique : on part de l'expression donnée, on repère dans le tableau~\ref{ch1:tab:lois} une loi applicable, on la remplace, et l'on recommence jusqu'à ce qu'aucune loi ne s'applique davantage.

\begin{methode}[simplifier une expression d'ensembles]
\begin{enumerate}
  \item éliminer d'abord les complémentaires d'expressions à l'aide des lois de De Morgan ;
  \item distribuer $\cap$ sur $\cup$ (ou $\cup$ sur $\cap$) pour faire apparaître les paires $A \cap \overline{A} = \varnothing$ et $A \cup \overline{A} = U$ ;
  \item simplifier à l'aide de l'idempotence, de l'absorption et des éléments neutre et absorbant ;
  \item contrôler le résultat en comparant les zones du diagramme de Venn ou en testant sur un petit exemple numérique.
\end{enumerate}
\end{methode}

\begin{exemple}[une simplification pas à pas]
Simplifions $E = (A \cap B) \cup (A \cap \overline{B}) \cup (\overline{A} \cap B)$.
\begin{align*}
E &= \bigl(A \cap (B \cup \overline{B})\bigr) \cup (\overline{A} \cap B)
   && \text{distributivité de } \cap \text{ sur } \cup\\
  &= (A \cap U) \cup (\overline{A} \cap B)
   && \text{car } B \cup \overline{B} = U\\
  &= A \cup (\overline{A} \cap B)
   && \text{neutre de } \cap\\
  &= (A \cup \overline{A}) \cap (A \cup B)
   && \text{distributivité de } \cup \text{ sur } \cap\\
  &= U \cap (A \cup B)
   && \text{car } A \cup \overline{A} = U\\
  &= A \cup B .
   && \text{neutre de } \cup
\end{align*}
L'expression de départ, qui mêle six opérations, vaut donc simplement $A \cup B$ : un étudiant suit programmation \emph{ou} réseaux, sans autre condition.
\end{exemple}

Le même ensemble se simplifie par plusieurs chemins ; l'ordre des lois n'est pas imposé, mais descendre d'abord les complémentaires puis distribuer donne presque toujours le chemin le plus court. En cas d'impasse, repartir de l'expression initiale : les erreurs de simplification se propagent, et un nouveau départ coûte moins que la recherche d'une faute de signe au milieu du calcul.

\begin{exemple}[passer par les lois de De Morgan]
Simplifions $F = \overline{\overline{A} \cup \overline{B}} \cup (A \cap B)$. La première parenthèse se traite par De Morgan :
\[
\overline{\overline{A} \cup \overline{B}} = \overline{\overline{A}} \cap \overline{\overline{B}} = A \cap B,
\]
en appliquant l'involutivité du complémentaire. L'expression devient $F = (A \cap B) \cup (A \cap B)$, que l'idempotence de $\cup$ réduit à $F = A \cap B$. Le contrôle par Venn est immédiat : la zone extérieure à « $A$ seul $\cup$ $B$ seul » est exactement le lentille central $A \cap B$.
\end{exemple}

\begin{attention}
Une simplification ne se termine que par un contrôle. Le réflexe le plus sûr est le test numérique : affecter aux ensembles $A$ et $B$ des petits ensembles concrets (deux sous-ensembles de $\{1, 2, 3, 4\}$ suffisent), puis calculer l'expression de départ et le résultat proposé avec les deux. Si les deux calculs diffèrent, la simplification est fausse ; s'ils coïncident, elle est très probablement juste, la méthode algébrique garantissant le reste. Sauter ce contrôle revient à confondre une réécriture élégante avec une réécriture correcte.
\end{attention}

\begin{aretirer}
Un ensemble se décrit en extension (liste des éléments) ou en compréhension (propriété d'appartenance) ; son cardinal $|A|$ compte ses éléments. Les trois opérations (union $A \cup B$, intersection $A \cap B$, complémentaire $\overline{A}$) se lisent sur le diagramme de Venn et s'ordonnent par les lois du tableau~\ref{ch1:tab:lois}, dont les deux lois de De Morgan sont les plus rentables : $\overline{A \cup B} = \overline{A} \cap \overline{B}$ et $\overline{A \cap B} = \overline{A} \cup \overline{B}$. Pour compter : $|A \cup B| = |A| + |B| - |A \cap B|$, un ensemble à $n$ éléments possède $2^{n}$ sous-ensembles, et $|A \times B| = |A| \times |B|$.
\end{aretirer}
